\section{Eigenschaften des adjungierten Operators $ S^\ast $}

\begin{sz}
	\begin{enumerate}[label=\text{\textbf{(\arabic*)}}]
		\item 
		Sei $ E $ ein abgeschlossener Teilraum von $ \H^2 $ mit $ E \neq \H^2 $ und $ S^\ast E \subset E $.
		Dann gilt 
		\begin{align*}
		E = \H^2 \ominus \Theta \H^2
		\end{align*}
		mit einer inneren Funktion $ \Theta $.
		Außerdem ist jeder Unterraum der Form $ \H^2 \ominus \Theta \H^2 $
		$ S^\ast $-  invariant.	
		
		\item
		Wenn $ E $ ein $ S $-und $ S^\ast $-invarianter Unterraum von $ \H^2 $ ist, dann gilt entweder $ E = {0} $ oder $ E = \H^2 $.
	\end{enumerate}
\end{sz}

\begin{proof}
	\begin{enumerate}[label=\text{\textbf{(\arabic*)}}]
		\item 
		Mit $ S^\ast E \subset E $ gilt
		\begin{align*}
		\langle x , S^\ast y \rangle =
		\langle S x ,  y \rangle = 0, \quad x \in E^\bot, y \in E,
		\end{align*}
		wodurch wir sofort
		\begin{align*}
		S^\ast E \subset E \Leftrightarrow S E^\bot \subset E^\bot
		\end{align*}
		erhalten. Mit Beurling-Helson erhalten wir dann $ E^\bot = \Theta \H^2 $, womit die Aussage gilt.
		
		\item
		Angenommen $ \lbrace 0 \rbrace \neq E \neq \H^2 $.
		Wegen $ S E \subset E  $ und $ S^\ast E \subset E $, gelten 
		\begin{align*}
		E = \Theta \H^2 = \H^2 \ominus \Theta^\prime \H^2
		\end{align*}
		für zwei innere Funktionen $ \Theta, \Theta^\prime$.
		Dann gilt 
		\begin{align*}
		\langle \Theta f , \Theta^\prime g \rangle   
		=
		\int \limits_{T} \Theta f \overline{\Theta^\prime g}  \td{\lambda_N} = 0
		\end{align*}
		für $ f, g \in \H^2 $. Für $ \langle \Theta \Theta^\prime, \Theta^\prime \Theta \rangle = 1 $ ist dies ein Widerspruch.		
		\end{enumerate}
\end{proof}

\begin{genericthm}{Folgerung}
	Sei $ \mathcal{F}$  eine Familie innerer Funktionen und 
	$ K_\Theta := \H^2 \ominus \Theta \H^2$ für eine innere Funktion $ \Theta $.
	Dann gelten:
	\begin{enumerate}[label=\text{\textbf{(\arabic*)}}]
		\item $ K_{\Theta_1} \subset K_{\Theta_2}  \ \Leftrightarrow \ \Theta_1 \ \text{teilt} \ \Theta_2$
		
		\item
		$ \bigvee_{\Theta \in \mathcal{F}} K_\Theta  = K_{\Theta^\prime}$ mit $ \Theta^\prime = \kgv(\mathcal{F})  $
		
		\item
		$ \bigcap_{\Theta \in \mathcal{F}} K_\Theta = K_{\Theta^\prime} $
		mit $ \Theta^\prime = \ggT(\mathcal{F}) $.
	\end{enumerate}
	Für den Fall $ \bigvee_{\Theta \in \mathcal{F}} K_\Theta  = \H^2 $ 
	und $ \bigcap_{\Theta \in \mathcal{F}} K_\Theta = \lbrace 0 \rbrace $ sagen wir $ \kgv( \mathcal{ F}) = 0 $ und $ \ggT(\mathcal{F}) = 1 $.
\end{genericthm}

\begin{proof}
	\begin{enumerate}[label=\text{\textbf{(\arabic*)}}]
		\item
		Durch
		\begin{align*}
		K_{\Theta_1} =  (\Theta_1 \H^2)^\bot  \subset (\Theta_2 \H^2)^\bot = K_{\Theta_2}
		\ \Leftrightarrow \ 
		\Theta_2 \H^2 \subset \Theta_1 \H^2
		\ \Leftrightarrow \
		\Theta_1 \ \text{teilt} \ \Theta_2
		\end{align*}
		erhalten wir die Aussage sofort.
		
		\item
		Der Raum 
		\begin{align*}
		\bigvee_{\Theta \in \mathcal{F}} K_\Theta
		\end{align*}
		ist $ S^\ast $-invariant.
		Damit existiert eine innere Funktion $ \Theta^\prime $ mit
		\begin{align*}
		\bigvee_{\Theta \in \mathcal{F}} K_\Theta = K_{\Theta^\prime}
		\end{align*}
		und es gilt
		\begin{align*}
		K_\Theta \subset K_{\Theta^\prime}
		\end{align*}
		für alle $ \Theta \in \mathcal{F} $.
		Damit gilt auch $ \Theta^\prime = \kgv(\mathcal{F}) $.
		
		\item
		Wegen der $ S^\ast $-Invarianz gilt auch hier 
		\begin{align*}
		\bigcap_{\Theta \in \mathcal{F}} K_\Theta = K_{\Theta^\prime}
		\end{align*}
		für eine innere Funktion $ \Theta^\prime $.
		Aus 
		\begin{align*}
		K_{\Theta^\prime} \subset K_{\Theta}
		\end{align*}
		für $ \Theta \in \mathcal{F} $ folgt direkt $ \Theta^\prime = \ggT(\mathcal{F}) $.
	\end{enumerate}
\end{proof}